\documentclass[10pt,a4paper]{amsart}
\usepackage[margin=19mm]{geometry}
\usepackage{amsmath,amssymb,mathtools}
\usepackage[scr=rsfs,cal=euler]{mathalfa}
\usepackage{xfrac}
\usepackage[unicode,hidelinks,bookmarks=false]{hyperref}
\newcommand{\ie}{i.e.\ }
\newcommand{\cf}{cf.\ }
\newcommand{\resp}{respectively\ }
\newcommand{\Subsection}{\subsection}
\providecommand{\nobreakdash}{\nobreak\hbox{-}}
\setlength{\emergencystretch}{2em}
\multlinegap0pt
\usepackage{amssymb}
\usepackage{xcolor}
\usepackage[normalem]{ulem}
\newtheorem{lemma}{Lemma}
\newtheorem{proposition}{Proposition}
\newcommand{\bZ}{\mathbb Z}
\def\ep{{\varepsilon}}
\title[Author proof 110147]{Zero distribution of power series and binary correlation of coefficients}
\author{Jacques Benatar, Alexander Borichev, Mikhail Sodin}
\date{Source: arXiv:2104.04812v1 (CC BY 4.0)}
\newcommand{\bC}{\mathbb C}
\newcommand{\bR}{\mathbb R}
\allowdisplaybreaks
\makeatletter
\newlabel{Lemma-subharm}{{\textup{source line 685}}{}}
\newlabel{LemmaL-W-V}{{\textup{source line 1194}}{}}
\newlabel{thm:main2}{{\textup{source line 2628}}{}}
\makeatother
\newtheorem{definition}{Definition}
\begin{document}

\maketitle
\noindent\textbf{Reviewed difficulty: H3. Legacy noncounting; count 0.} The H3 burden includes quantitative autocorrelation and no-gap spectral-mass hypotheses, the complete dedicated Gaussian-kernel chain at 1546--1823 (with its regularity estimates at 1042--1113), and the $I+II+III$ lower-bound construction at 2454--2621, allowing singular spectral measures. These are joint research arguments in this unit. Source note: the positive exponent printed at 2464 is retained literally. The unique negative-Gaussian definition at 1885--1886 and the same Fourier identity at 1910--1928, explicitly reused at 2455 and 2468, identify the omitted minus sign. This resolved source notation issue does not create an active hold or a replacement author proof.\par\medskip

\noindent\emph{Source and licence.} Zero distribution of power series and binary correlation of coefficients. Version arXiv:2104.04812v1 (CC BY 4.0), distributed by arXiv under the Creative Commons Attribution 4.0 International licence, \url{http://creativecommons.org/licenses/by/4.0/}, as recorded in the arXiv licence metadata for this exact version and shown on the abstract page \url{https://arxiv.org/abs/2104.04812v1}. Full licence notice: \path{licenses/arxiv-2104.04812-CC-BY-4.0.txt}.\par\smallskip


\noindent\textit{Editorial scope.} Selected author statement at source0.tex 2427--2452 and complete author proof at 2454--2621. Necessary author inputs and any dedicated auxiliary chain are displayed below. Original mathematical text, hypotheses, formulas and source version are retained.\par\smallskip
\section*{Author context (source 517--541)}
\begin{itemize}
\item $e(t)=e^{2\pi {\rm i}t}$.
\item $F_\xi(z) = \displaystyle \sum_{n\ge 0} \xi(n) a(n) z^n$ the entire function.
\item $a(n) = \displaystyle \exp\Bigl[\, -\int_0^n \varphi \,\Bigr]$ ``smooth coefficients''.
\item $\psi = \varphi^{-1}$ the inverse of $\varphi$.
\item $\mu = \displaystyle \max_{t\ge 0}\, \exp\Bigl[\,
t\log R - \int_0^t \varphi \,\Bigr]$ the smoothed maximal term.
\item $\nu = \psi\circ\log$ the smoothed central index.
\item $\sigma = \psi'\circ\log$.
\item $\gamma=(2\pi)^{-1}\Delta\log\mu(|z|) =
\sigma (r)r^{-1}\, {\rm d}r \otimes {\rm d}\theta$, $z=r e(\theta)$, the reference measure.
\item $\rho$ the radial gauge on $\bC$.
\item $d_\rho (z_1, z_2) = \displaystyle \inf_\ell\, \int_\ell \frac{|{\rm d}w|}{\rho(w)}$,
the infimum is taken over all $C^1$ curves $\ell$ connecting $z_1$ with $z_2$, the distance in $\bC$.
\item $U_{+\tau} = \{z\in\bC\colon d_\rho(z, U)<\tau\}$ the $\tau$-neighbourhood of the set
$U\subset \bC$.
\item $W_R(\theta) = \displaystyle
\sum_{|n-\nu|\le N} \xi (n)e(n\theta) e^{-(n-\nu)^2/(2\sigma)}$,
$N=A\sqrt{\sigma\log\sigma}$ with sufficiently large positive $A$, the Weyl-type exponential sum.
\item $A\lesssim B$ means that $A\le CB$ with a positive constant $C$,
$A \gtrsim B$ means $A\ge cB$ with a constant $c>0$, and $A\simeq B$ means that $A\lesssim B$ and
$A\gtrsim B$ simultaneously.
\item The sign $\ll$ means ``sufficiently smaller than'' and $\gg$ means ``sufficiently large than''. For instance, the assumption `` given $A$ and $B$ such that $A\ll B$'' means that there exists
    $c\in (0, 1)$ such that the corresponding conclusion holds for every positive $A$ and $B$ satisfying $A\le cB$.
\end{itemize}

\section*{Author context (source 1042--1113)}
\begin{definition}[$\Delta$-regularity]
Let $\Delta\colon (0, \infty)\to (0, \infty)$ be a non-decreasing function satisfying
$\Delta (2s) \simeq \Delta (s)$ and $ \Delta(s) = o( \sqrt{s}\,(\log s)^{-3/2})$ as $s\to\infty$.
The function $\varphi$ will be called $\Delta$-regular,
if it is a non-negative, increasing, concave $C^2$-function on $[0, \infty)$, satisfying
$$
\varphi (0) = 0,
\quad \lim_{t\to\infty} \varphi (t) = \infty, \quad \lim_{t\to\infty} \varphi' (t)= 0\,,
$$
with
\begin{equation}\label{eq:phi-very-new}
|\varphi''| \le (\varphi')^2\, \Delta(1/\varphi')\,.
\end{equation}
\end{definition}

By $\psi=\varphi^{-1}$ we denote the inverse function to $\varphi$.
This is a convex function, which, together with its derivative, grows to $+\infty$.

We start with some simple estimates for $\Delta$-regular functions $\varphi$ and their inverses
$\psi$, which will be used throughout this work.

\begin{lemma}\label{Lemma-AuxEstimates-new}
Suppose that the function $\varphi$ is $\Delta$-regular.
Then,
\begin{itemize}
\item[{\rm (a)}] $\psi''\le \psi'\Delta(\psi')$;
\item[{\rm (b)}] $\varphi'(t)=(1+o(1))\varphi'(\tau)$, for
$\tau\to\infty$ and
$|t-\tau|\varphi'(\tau)\Delta(1/\varphi'(\tau))=o(1)$;
\item[{\rm (c)}] $\psi'(u) = (1+o(1))\psi'(v)$, for
$v\to\infty$ and $|u-v|\Delta(\psi'(v))=o(1)$;
\item[{\rm (d)}] $\varphi'(t)\Delta(1/\varphi'(t)) \gtrsim 1/t$;
\item[{\rm (e)}] $\psi'/\Delta(\psi') \lesssim \psi$.
\end{itemize}
\end{lemma}

\begin{proof}
Estimate (a) follows from~\eqref{eq:phi-very-new} combined with the formulas
\[
\psi' = \frac1{\varphi'(\psi)}\,,
\qquad \psi''=\frac{|\varphi''(\psi)|}{(\varphi'(\psi))^3}\,.
\]

Suppose that estimate (b) does not hold, i.e., that there exists a function $\ep(\tau)\to 0$
as $\tau\to\infty$, such that
\[
\limsup_{\tau\to\infty}\, \max\{|\varphi'(t)/\varphi'(\tau)-1|\colon |t-\tau|\varphi'(\tau)\Delta(1/\varphi'(\tau)) \le \ep(\tau) \} > 0.
\]
Then, for some $\delta\in (0, \tfrac12)$, there exist an arbitrarily large $\tau$ and a
$\vartheta$ such that $|\varphi' (\vartheta)/\varphi'(\tau)-1|=\delta$, while
$|\varphi' (t)/\varphi'(\tau)-1|<\delta$ everywhere on the open interval $I$ with endpoints $\tau$ and $\vartheta$ and with $ |I|\varphi'(\tau)\Delta(1/\varphi'(\tau)) \le \ep(\tau) $.
Then,
\[
\delta \lesssim \Bigl| \log\,\frac{\varphi'(\vartheta)}{\varphi'(\tau)}\, \Bigr|
\le \int_I\, \Bigl|\, \frac{\varphi''}{\varphi'}\, \Bigr| \le |I|\cdot \max_I \Bigl|\, \frac{\varphi''}{\varphi'}\, \Bigr| \le |I|\cdot \max_I \varphi' \Delta(1/\varphi').
\]
Since $\delta$ was chosen less than $\tfrac12$,
everywhere on $I$ we have $ \varphi' \le 2\varphi'(\tau) $. Hence,
$ \max_I \varphi' \Delta(1/\varphi') \lesssim \varphi'(\tau) \Delta(1/\varphi'(\tau)) $, and therefore,
\[
\delta \lesssim |I| \cdot \varphi'(\tau) \Delta(1/\varphi'(\tau)) \lesssim \ep(\tau)\,,
\]
arriving at a contradiction as $\tau\to\infty$, which proves (b).
The proof of estimate (c) follows the same pattern, so we skip it.

Estimate (d) easily follows from (b). Indeed, assume that (d) does not hold.
That is, there is a sequence $\tau_j\uparrow\infty$ such that
$\tau_j \varphi'(\tau_j)\Delta(1/\varphi'(\tau_j))\to 0$. Consider the intervals
$[0, \tau_j]$. By (b), $|\varphi'(+0)/\varphi'(\tau_j)-1|\to 0$. Since $\varphi'(\tau_j)\to 0$,
we get $\varphi'(+0) = 0$, which is impossible since $\varphi'(\tau)$ monotonically decreases to $0$
as $\tau$ grows. At last, estimate (e) is just a restatement of (d).
\end{proof}

\section*{Author context (source 1546--1823)}
\subsection{Auxiliary estimates}
Here, we collect estimates of the function \newline\noindent $V_R(k; h)$, which will be defined
momentarily. These estimates will be used in the proofs of Proposition~\ref{Lemma5}
and Proposition~\ref{Lemma5a}.

Throughout this section we assume that the function $\varphi$ is $\Delta$-regular.
As above, $\nu=\nu(R)$ and $\sigma=\sigma(R)$. Let $\beta=\beta(R)$ be a small parameter
satisfying
\[
\sigma^{-1/2} = o(\beta), \quad
\beta\Delta(\sigma) = o(1)\,,
\]
as $R\to\infty$. Set
\begin{align*}
V_R(k; h) &= \int_R^{R(1+\beta)}
\exp\Bigl[ -\frac{(k-\nu(s))^2 + (k+h-\nu(s))^2}{2\sigma(s)}\, \Bigr]\, {\rm d}\nu(s) \\
&=
\int_{\nu(R)}^{\nu(R(1+\beta))}
\exp\Bigl[ -\frac{(k-t)^2 + (k+h-t)^2}{2\sigma(\nu^{-1}(t))}\, \Bigr]\, {\rm d}t\,.
\end{align*}
A simple and useful observation is that the function $\sigma = \psi'\circ\log $ stays
approximately constant on the integration interval $[R, R(1+\beta)]$:

\begin{lemma}\label{Lemma5A}
We have
\begin{itemize}
\item[{\rm (i)}]
$\sigma (s) = (1+o(1))\sigma $ everywhere on $[R, R(1+\beta)]$;
\item[{\rm (ii)}]
$\nu(R(1+\beta))-\nu (R) = (1+o(1))\beta\sigma$.
\end{itemize}
\end{lemma}

\begin{proof}
Since $|\log(R(1+\beta))-\log R \,| \le \beta = o(1/\Delta(\sigma))$,
relation (i) readily follows from Lemma~\ref{Lemma-AuxEstimates-new}(c).
To prove (ii), we recall that $\nu'(r) = \sigma (r)/r$. Hence,
(ii) follows from (i).
\end{proof}

The next lemma provides us with crude upper and lower bounds on $V_R$.

\begin{lemma}\label{Lemma5C}
We have
\begin{itemize}
\item[{\rm (i)}] $V_R(k; h)\lesssim \sqrt{\sigma}$;
\item[{\rm (ii)}] $V_R(k; 0) \gtrsim \sqrt{\sigma}$,
provided that  $\nu(R)\le k \le \nu(R(1+\beta))$.
\end{itemize}
\end{lemma}

\begin{proof}
By Lemma~\ref{Lemma5A}(i),
\begin{align*}
V_R(k; h) &\le \int_{\nu(R)}^{\nu(R(1+\beta))}
\exp\bigl[ - (1+o(1))\,(k-t)^2/(2\sigma)\bigr]\, {\rm d}t \\
&<\int_\bR \exp\bigl[ -(1+o(1))u^2/{2\sigma} \bigr]\, {\rm d}u \\
&\lesssim\sqrt{\sigma}\,,
\end{align*}
proving (i).

The proof of (ii) is also straightforward. By Lemma~\ref{Lemma5A}(i),
\begin{align*}
V_R(k; 0) &=
\int_{\nu(R)}^{\nu(R(1+\beta))} \exp\bigl[ - (1+o(1))(k-t)^2/(2\sigma) \bigr]\, {\rm d}t \\
&=
\int_{\nu(R)-k}^{\nu(R(1+\beta))-k} \exp\bigl[ - (1+o(1))u^2/(2\sigma) \bigr]\, {\rm d}u\,.
\end{align*}
Observe that either $\nu (R) - k \le -\tfrac12\,L$, or $\nu (R(1+\beta))-k \ge \tfrac12\, L$,
where $L=\nu (R(1+\beta))-\nu (R)$. By Lemma~\ref{Lemma5A}(ii), $L=(1+o(1))\beta\sigma$, which is
$\gg \sqrt{\sigma}$. Thus,
\[
\int_0^{\frac12\, L} \exp\bigl[ - (1+o(1))u^2/(2\sigma) \bigr]\, {\rm d}u \gtrsim \sqrt{\sigma}\,,
\]
proving the lemma.
\end{proof}

The next lemma gives us more accurate upper bounds.

\begin{lemma}\label{Lemma5B}\mbox{}
\begin{itemize}
\item[{\rm (i)}] $V_R(k; h) \lesssim e^{-ch^2/\sigma}\sqrt{\sigma}$;
\item[{\rm (ii)}]
\[
V_R(k; h) \lesssim
\begin{cases}
e^{-c(\nu(R)-k)^2/\sigma}\sqrt{\sigma}, \ &k<\nu(R), \\
e^{-c(k-\nu(R(1+\beta)))^2/\sigma}\sqrt{\sigma}, \ &k>\nu(R(1+\beta))\,.
\end{cases}
\]
\end{itemize}
\end{lemma}

\begin{proof}
We have
\begin{align*}
V_R(k; h) &=
\int_{\nu (R)-k-\frac{h}2}^{\nu(R(1+\beta))-k-\frac{h}2}
\exp\Bigl[ -\frac{1+o(1)}{2\sigma}\,
\bigl( (u-\tfrac12 h)^2 + (u+\tfrac12 h)^2 \bigr)
\Bigr]\, {\rm d}u \\
&=
\int_{\nu (R)-k-\frac{h}2}^{\nu(R(1+\beta))-k-\frac{h}2}
\exp\bigl[ -(1+o(1))(u^2 + \tfrac14 h^2)/\sigma \bigr]\, {\rm d}u \\
&<
e^{-(1+o(1))h^2/(4\sigma)}\,
\int_{\bR} e^{-(1+o(1))u^2/\sigma}\, {\rm d}u \\
&\lesssim e^{-ch^2/\sigma}\, \sqrt{\sigma}\,.
\end{align*}
To prove the second estimate, we note that, for $k<\nu(R)$, we have
\begin{align*}
V_R(k; h) &<
\int_{\nu(R)}^{\nu(R(1+\beta))}
\exp\Bigl[ -\frac{1+o(1)}{2\sigma}\, (t-k)^2 \Bigr]\, {\rm d}t \\
&<
\int_{\nu(R)-k}^\infty e^{-(1+o(1))u^2/(2\sigma)}\, {\rm d}u \\
&< e^{-(1+o(1))(\nu(R)-k)^2/(2\sigma)}\sqrt{\sigma}\,.
\end{align*}
The case $k>\nu(R(1+\beta))$ is very similar and we skip it.
\end{proof}

The next lemma estimates the oscillation of the function $k\mapsto V_R(k; h)$ along $\mathbb Z$.

\begin{lemma}\label{Lemma5D}
$ \displaystyle \sum_{k\in\bZ}
|V_R(k; h) - V_R(k-1; h)| \lesssim \sqrt{\sigma} $.
\end{lemma}

\begin{proof}
Set $m_1 = \nu(R) - A\sqrt{\sigma\log\sigma}$, $m_2 = \nu(R(1+\beta)) +A\sqrt{\sigma\log\sigma}$
with sufficiently large $A$. By Lemma~\ref{Lemma5B}(ii), the sums over $k<m_1$ and
$k>m_2$
are $o(1)$ if the constant $A$ was chosen big enough. Hence, we need to show that
\[ \sum_{m_1 \le k \le m_2} |V_R(k; h) - V_R(k-1; h)| \lesssim \sqrt{\sigma}\,. \]

We will represent the function $V_R(k; h)$ as a difference of two increasing
functions in $k$. Recalling that $\sigma (\nu^{-1}(t)) = 1/\varphi'(t))$,
we get
\begin{align*}
V_R(k; h) &= \int_{\nu(R)}^{\nu(R(1+\beta))}
\exp\Bigl[\,
- \frac{(k-t)^2+(k+h-t)^2}{2\sigma(\nu^{-1}(t))}
\,\Bigr]\, {\rm d}t \\
&= \int_{\nu(R)}^{\nu(R(1+\beta))}
\exp\bigl[\,
- \tfrac12\, \varphi'(t)((k-t)^2+(k+h-t)^2)
\,\bigr]\, {\rm d}t \\
&= \int_{\nu(R)-k-\frac12 h}^{\nu(R(1+\beta))-k-\frac12 h}
\exp\bigl[\,
- \tfrac12\, \varphi'(x+k+\tfrac12\, h)(x^2 + \tfrac14\, h^2)
\,\bigr]\, {\rm d}x \\
&= \Bigl(\, \int_{\nu(R)-k-\frac12 h}^{\nu(R(1+\beta))-m_1 -\frac{h}2}
- \int_{\nu(R(1+\beta))-k-\frac12 h}^{\nu(R(1+\beta))-m_1 -\frac{h}2} \,\Bigr)\\
&\qquad\qquad\qquad\qquad\qquad\exp\bigl[
- \tfrac12\, \varphi'(x+k+\tfrac12\, h)(x^2 + \tfrac14\, h^2)
\,\bigr]\, {\rm d}x \\
&= V_{1, R}(k; h) - V_{2, R}(k; h).
\end{align*}

First, observe that both functions $V_{1, R}$ and $V_{2, R}$ have uniform upper bounds
\begin{equation}\label{eq:UB-V_1-V_2}
V_{1, R}, V_{2, R} \lesssim \sqrt{\sigma}\,.
\end{equation}
Indeed, in the integration range, $\nu(R) \le x+k+\tfrac12\, h \le \nu(R(1+\beta)) + (m_2-m_1)
= \nu + O(\beta\sigma)$. Since $\beta\sigma\Delta(\sigma) \varphi'(\nu) =
\beta \Delta(\sigma) = o(1)$, by Lemma~\ref{Lemma-AuxEstimates-new}(b), in this range,
$\varphi'(x+k+\tfrac12\, h) = (1+o(1))\varphi'(\nu) = (1+o(1))\sigma^{-1}$,
which yields~\eqref{eq:UB-V_1-V_2}.

Next, we observe that since the function $t\mapsto \varphi'(t)$ decreases,
the functions $V_{1, R}$ and $V_{2, R}$ are increasing functions in $k$.

The rest is straightforward:
\begin{align*}
 \sum_{m_1 \le k \le m_2}
|V_R(k; h) &- V_R(k-1; h)| \\
& 
\le
\sum_{1\le j\le 2}\,\sum_{m_1 \le k \le m_2}
|V_{j, R}(k; h) - V_{j, R}(k-1; h)|
\\
& 
= \sum_{1\le j\le 2}\,\sum_{m_1 \le k \le m_2}
\bigl[ V_{j, R}(k-1; h) - V_{j, R}(k; h) \bigr]
\\
&   < \sum_{1\le j\le 2}V_{j, R}(m_1; h) 
\\
&  \lesssim \sqrt{\sigma}\,,
\end{align*}
proving the lemma.
\end{proof}

Set
\[
V_R(h) = \sum_{k\in\bZ} V_R(k; h)\,.
\]

\begin{lemma}\label{Lemma10a} We have
\[
V_R(0) \simeq \beta \sigma^{3/2}\,.
\]
\end{lemma}

\begin{proof}
To prove the lower bound, we write
\[
V_R(0) \ge \sum_{\nu(R)\le k \le \nu(R(1+\beta))}  V_R(k; 0)
\]
and note that by Lemma~\ref{Lemma5A}(ii) and Lemma~\ref{Lemma5C},
the RHS is
\[
\simeq \sqrt{\sigma} \bigl[ \nu(R(1+\beta)) - \nu (R) \bigr] \simeq \beta \sigma^{3/2}.
\]
To prove the upper bound, we split the sum into three parts:
\[
V_R(0) = \Bigl( \sum_{k<\nu (R)}\ + \sum_{\nu(R)\le k \le \nu(R(1+\beta))}
+ \sum_{k>\nu(R(1+\beta))}\, \Bigr) V_R(k; 0)\,.
\]
By Lemma~\ref{Lemma5B}(ii), the first and the third sums are $O(\sigma)=o(\beta\sigma^{3/2})$,
while, by the above, the middle sum is $\simeq \beta\sigma^{3/2}$.
\end{proof}

\begin{lemma}\label{Lemma10b} For $ h^2 \lesssim \sigma $, we have
\[
|V_R(h) - V_R(0)| \lesssim (1+ h^2\beta)\sqrt{\sigma}\,.
\]
\end{lemma}

\begin{proof}
First, we observe that
\[
\Bigl| V_R(h) - \int_{\bR} V_R(\kappa; h)\, {\rm d}\kappa \Bigr|
\lesssim \sqrt{\sigma}\,.
\]
Indeed, as in the proof of Lemma~\ref{Lemma5D}, we set
$m_1 = \nu(R) - A\sqrt{\sigma\log\sigma}$, $m_2 = \nu(R(1+\beta)) +A\sqrt{\sigma\log\sigma}$
with sufficiently large $A$. Then, by Lemma~\ref{Lemma5B}(ii),
\[
V_R(h) - \int_{\bR} V_R(\kappa; h)\, {\rm d}\kappa
= \sum_{m_1 \le k \le m_2} \Bigl[ V_R(k; h) -
\int_k^{k+1} V_R(\kappa; h)\, {\rm d}\kappa \Bigr] + o(1)\,,
\]
provided that the constant $A$ was chosen sufficiently big. It remains to recall that,
as we have shown in the proof of Lemma~\ref{Lemma5D}, the total variation of the function
$k\mapsto V_R(k; h)$ on $[m_1, m_2]$ is $\lesssim\sqrt{\sigma}$.

Thus, we need to bound the difference of the integrals
\[
\Bigl| \int_{\bR} V_R(\kappa; h)\, {\rm d}\kappa -
\int_{\bR} V_R(\kappa; 0)\, {\rm d}\kappa\Bigr|
\le \int_{m_1}^{m_2} \bigl| V_R(\kappa; h) - V_R(\kappa-\tfrac12\, h; 0) \bigr|\, {\rm d}\kappa
+ o(1)\,.
\]
Estimating the integrand on the RHS, we get
\begin{align*}
\bigl|&V_R(\kappa; h) - V_R(\kappa-\tfrac12\, h; 0) \bigr| \\
&=
\Bigl|
\int_{\nu(R)-\kappa-\frac12\, h}^{\nu (R(1+\beta))-\kappa-\frac12\, h}
\bigl[ \exp\bigl( -\varphi'(u+\kappa+\tfrac12\, h)(u^2 + \tfrac14\, h^2)\,\bigr)\\
&\qquad\qquad\qquad\qquad\qquad\qquad\qquad\qquad -\exp\bigl( -\varphi'(u+\kappa+\tfrac12\, h)u^2 \bigr]\, {\rm d}u
\Bigr| \\
&\le \int_{\nu(R)-\kappa-\frac12\, h}^{\nu (R(1+\beta))-\kappa-\frac12\, h}
\exp\bigl( - (1+o(1))u^2/\sigma \bigr)\\
&\qquad\qquad\qquad\qquad\qquad\qquad\qquad\qquad\times\bigl| \exp\bigl( -(1+o(1))h^2/(4\sigma) \bigr) - 1\bigr|\, {\rm d}u
\\
&\lesssim (h^2/\sigma) \cdot \sqrt{\sigma} = h^2/\sqrt{\sigma}\,,
\end{align*}
Thus,
\begin{multline*}
\Bigl| \int_{\bR} V_R(\kappa; h)\, {\rm d}\kappa -
\int_{\bR} V_R(\kappa; 0)\, {\rm d}\kappa\Bigr|
\lesssim (m_2-m_1)h^2/\sqrt{\sigma} + o(1) \\
\lesssim (\beta\sigma + \sqrt{\sigma\log\sigma})h^2/\sqrt{\sigma}
\lesssim h^2\beta\sqrt{\sigma}\,,
\end{multline*}
completing the proof.
\end{proof}

\section*{Author context (source 1847--1870)}
\begin{proposition}\label{Lemma5}
Suppose that
\begin{equation}
\label{eq:condition1}
\sum_{\nu \le k \le \nu+\frac12 \beta\sigma}\, |\xi(k)|^2 \gtrsim \beta\sigma\,,
\end{equation}
and that, for some $p>1$,
\begin{equation}
\label{eq:condition2}
\sum_{h \ge 1} (1+\beta h)^{-p} S^*(M_1, M_2; h) \ll \beta\sigma,
\end{equation}
with
\[
M_1 = \nu (R) - \beta\sigma, \quad M_2 =  \nu(R(1+\beta)) + \beta\sigma\,.
\]
Then, for every $\vartheta\in [-\frac12, \frac12]$, there exist
\[
R'\in [R, R(1+\beta)], \quad \theta'\in (\vartheta-\beta, \vartheta+\beta),
\]
such that
\[
|W_{R'}(\theta')| \gtrsim \sigma^{1/4}\,.
\]
\end{proposition}

\section*{Author context (source 1882--1929)}
\noindent\textit{Author proof excerpt 1882--1929, showing the previously defined Gaussian and Fourier identity. The earlier proof continues outside this excerpt.}
To simplify our notation, we extend the sequence $\xi$, letting $\xi$ equal $0$
on negative integers, and set
\[
\widetilde{W}_R(\theta)
= \sum_{k\in\bZ} \xi(k)e(k\theta)\exp\Bigl[ - \frac{(k-\nu)^2}{2\sigma} \Bigr].
\]
It's easy to see that, for $N = A\,\sqrt{\sigma\log\sigma}$ with sufficiently large positive
$A$ (used in the definition of the sum $W_R$), we have
\[
\sum_{|k-\nu|\ge N} \exp\Bigl[ - \frac{(k-\nu)^2}{2\sigma} \Bigr] \lesssim
\int_N^\infty e^{-x^2/(2\sigma)}\, {\rm d}x = o(1),
\]
so that, in order to prove Proposition~\ref{Lemma5}, it will be enough to estimate the sum
$\widetilde{W}_R$, rather than $W_R$, from below.

We fix a non-negative function $g\in C_0^\infty(\bR)$, such that $\operatorname{supp}(g)\subset (-\tfrac12, \tfrac12)$, and $\displaystyle \int_\bR g = 1$, fix $\vartheta\in [-\tfrac12, \tfrac12]$,
and set
\[
\widetilde{X} =
\int_R^{R(1+\beta)}\, \int_{-1/2}^{1/2} |\widetilde{W}_s(\theta)|^2\,
g(\beta^{-1}(\vartheta-\theta))\, {\rm d}\theta\, {\rm d}\nu (s)\,.
\]
By Lemma~\ref{Lemma5A}(ii), $\nu(R(1+\beta))-\nu(R) = (1+o(1))\beta\sigma$, so, to prove Proposition~\ref{Lemma5}, 
we need to show that $\widetilde{X} \gtrsim \beta^2\sigma^{3/2}$.

First, we rewrite $\widetilde X$ as a Fourier series
\begin{align*}
\widetilde{X} &=
\int_R^{R(1+\beta)}\, \underset{k_1, k_2\in\bZ}{\sum\, \sum}\, \xi(k_1)\bar\xi(k_2)\,
\Bigl[
\int_{-1/2}^{1/2} e((k_1-k_2)\theta) g(\beta^{-1}(\vartheta-\theta))\, {\rm d}\theta
\Bigr] \\
& \qquad \qquad \qquad \qquad \qquad
\cdot \exp\Bigl[ -\frac{(k_1-\nu(s))^2 + (k_2-\nu(s))^2}{2\sigma(s)}\, \Bigr]\, {\rm d}\nu(s) \\
&= \beta\, \int_R^{R(1+\beta)}\, \underset{k_1, k_2\in\bZ}{\sum\, \sum}\,
e((k_1-k_2)\vartheta) \widehat{g}(\beta(k_2-k_1))\, \xi(k_1)\bar \xi(k_2) \\
& \qquad \qquad \qquad \qquad \qquad
\cdot \exp\Bigl[ -\frac{(k_1-\nu(s))^2 + (k_2-\nu(s))^2}{2\sigma(s)}\, \Bigr]\, {\rm d}\nu(s) \\
&= \beta\,
\sum_{h\in\bZ} \widehat{g}(\beta h) e(-h\vartheta)\,
\sum_{k\in\bZ} \xi(k)\bar\xi(k+h) V_R(k; h)\,,
\end{align*}
where
\begin{align*}
V_R(k; h) &= \int_R^{R(1+\beta)} \exp\Bigl[ -\frac{(k-\nu(s))^2 + (k+h-\nu(s))^2}{2\sigma(s)}\, \Bigr]\, {\rm d}\nu(s) \\
&=
\int_{\nu(R)}^{\nu(R(1+\beta))} \exp\Bigl[ -\frac{(k-t)^2 + (k+h-t)^2}{2\sigma(\nu^{-1}(t))}\, \Bigr]\, {\rm d}t \,.
\end{align*}

\section*{Author context (source 2388--2426)}
\section{Wiener sequences $\xi$ whose spectral measures have no gaps}
\label{sect:8}

Throughout this section we assume that
\begin{equation}\label{eq:C1}
\Bigl|\, \frac1X\, \sum_{0\le s < X} \xi(s)\bar\xi(s+h) - \widehat{\chi}(h) \Bigr|
\lesssim \ep_1(X; h), \qquad 0\le h \le H=H(X),
\end{equation}
with $X\mapsto \ep_1(X; h)$ decreasing to $0$ and $X\mapsto H(X)$ increasing to $\infty$,
as $X\to\infty$, and that
\begin{equation}\label{eq:C2}
\inf\bigl\{\chi(J)\colon
J\subset [-\tfrac12, \tfrac12\,] {\rm\ an\ interval}, |J|=\tfrac12\, \eta\bigr\}
\gtrsim \ep_2(\eta),
\end{equation}
with a positive non-decreasing function $\ep_2$.
The first condition quantifies the fact that $\chi$ is a spectral measure
of the Wiener sequence $\xi$, while the second condition is a quantitative version
of the statement that $\chi$ has no gaps in its support.

In Proposition~\ref{Lemma5a} we provide a set of conditions which will yield a lower
bound on the Weyl-type sum $W_R(\theta)$ on a sufficiently dense set of points
$(R, \theta)$. Then, we combine Proposition~\ref{Lemma5a} with Proposition~\ref{Lemma-subharm} and Proposition~\ref{LemmaL-W-V} and show (in Theorem~\ref{thm:main2}) that these
conditions guarantee equidistribution of zeroes of $F_\xi$ on appropriate local scales.
This set of conditions looks somewhat cumbersome, but then, to demonstrate how neatly it works,
we consider two instances of Wiener sequences $\xi$  with singular spectral measures
having no gaps in their support: the indicator-function of the square-free integers
and the Thue--Morse sequence.

\subsection{Another lower bound for Weyl-type sums}
\label{subsec:LB-general}

Denote by $\mathcal G$ the class of non-negative test-functions $g\in C_0^\infty (\bR)$ such that
$\operatorname{supp}(g)\subset \bigl(-\tfrac12, \tfrac12\, \bigr)$, $\displaystyle \int_\bR g = 1$, and
$g=1$ on $\bigl[-\tfrac14, \tfrac14\, \bigr]$.
As before, we assume that $\varphi$ is a $\Delta$-regular function, and
that $\nu = \psi\circ\log$, $\sigma=\psi'\circ\log$, where $\psi=\varphi^{-1}$ is the inverse
function.

\section*{Selected author statement}
\begin{proposition}\label{Lemma5a}
Let $R\gg 1$. Suppose that there exist $q>1$, $\beta=\beta(R)\ll 1/\Delta(\sigma)$,
and $g\in \mathcal G$, satisfying the following set of conditions:

\smallskip\noindent{\rm (a)}
$\beta^{-q} \le \min\bigl\{\sqrt{\sigma}, H(\tfrac12\, \nu)\bigr\}$;

\medskip\noindent{\rm (b)}
$\beta^{-(1+2q)} \ll \sigma\ep_2(\beta)$;

\medskip\noindent{\rm (c)}
$\displaystyle \nu \sum_{0\le h \le \beta^{-q}} \ep_1(\tfrac12\, \nu; h) \ll \beta\sigma\ep_2(\beta) $;

\medskip\noindent{\rm (d)}
$\displaystyle \sum_{h>\beta^{-q}} |\widehat{g}(\beta h)| \ll \ep_2(\beta) $.

\smallskip
Then, for every $\vartheta\in \bigl[-\tfrac12,  \tfrac12\, \bigr]$, there exist
\[
R'\in [R, R(1+\beta)],  \quad \theta'\in (\vartheta-\beta, \vartheta+\beta),
\]
such that
\[
|W_{R'}(\theta')| \gtrsim \sigma^{1/4}\, \sqrt{\ep_2 (\beta)}\,.
\]
\end{proposition}
\section*{Complete author proof}
\begin{proof}
As in the proof of Proposition~\ref{Lemma5}, we estimate from below the average
\[
\widetilde{X} =
\int_R^{R(1+\beta)}\, \int_{-1/2}^{1/2} |\widetilde{W}_s(\theta)|^2\,
g(\beta^{-1}(\vartheta-\theta))\, {\rm d}\theta\, {\rm d}\nu (s)\,,
\]
where
\[
\widetilde{W}_R(\theta) = \sum_{k\in\bZ} \xi(k)e(k\theta)
\exp\bigl[\,\frac{(k-\nu)^2}{2\sigma} \,\bigr]\,.
\]
By Lemma~\ref{Lemma5A}(ii), $\nu(R(1+\beta))-\nu(R) = (1+o(1))\beta\sigma$, so, to prove Proposition~\ref{Lemma5a} we need to show that
$\widetilde{X} \gtrsim \beta^2\sigma^{3/2}\ep_2(\beta)$.
As before, we rewrite $\widetilde X$ as a Fourier series
\[
\widetilde{X}
= \beta\,
\sum_{h\in\bZ} \widehat{g}(\beta h) e(-h\vartheta)\,
\sum_{k\in\bZ} \xi(k)\bar\xi(k+h) V_R(k; h)\,.
\]
Recalling the notation
$\displaystyle V_R(h) = \sum_{k\in\bZ} V_R(k; h)$, we split the RHS into three parts:
\begin{align*}
\widetilde{X}  &= \beta V_R(0)\, \sum_{h\in\bZ} \widehat{g}(\beta h)\widehat{\chi}(h) e(-h\vartheta) \\
&\quad + \beta\, \sum_{h\in\bZ} \widehat{g}(\beta h) e(-h\vartheta)\,
\sum_{k\in\bZ} \bigl( \xi(k)\bar\xi(k+h) - \widehat{\chi}(h) \bigr) V_R(k; h) \\
&\quad + \beta\, \sum_{h\in\bZ} \widehat{g}(\beta h) \widehat{\chi}(h) e(-h\vartheta)\,
(V_R(h) - V_R(0)) \\
&= I +II + III\,.
\end{align*}
We will show that $I\gtrsim \beta^2\sigma^{3/2}\ep_2(\beta)$, while
the terms $|II|$ and $|III|$ are  $\ll \beta^2\sigma^{3/2}\ep_2(\beta)$.

\medskip\noindent\underline{Lower bound of $I$}:
We set $g_\beta(\theta) = \beta^{-1}g(\beta^{-1}\theta)$, and denote by
$(\chi * g_\beta)'$ the density of the convolution of $\chi$ with
$g_\beta$. Then, \[ I=\beta V_R(0) (\chi * g_\beta)'(-\vartheta).\]
By Lemma~\ref{Lemma10a}, $V_R(0)\gtrsim \beta \sigma^{3/2}$. Since $g_\beta = \beta^{-1}$ on
$ \bigl[-\tfrac14\, \beta, \tfrac14\, \beta \bigr]$, we have
\[
(\chi * g_\beta)'(-\vartheta) \gtrsim
\chi\bigl[-\vartheta-\tfrac14\, \beta, -\vartheta + \tfrac14\, \beta \bigr]
\gtrsim \ep_2(\beta).
\]
Thus, $I \gtrsim \beta^2 \sigma^{3/2} \ep_2(\beta)$.

\medskip\noindent\underline{Upper bound of $II$}:
Recalling that $\beta\sigma \stackrel{(a)}\gg \sqrt{\sigma\log\sigma}$
and using Lemma~\ref{Lemma5B}(ii), we cut the sum in $k$, getting
\begin{align*}
|II| &\lesssim \beta\, \sum_{h\in\bZ} |\widehat{g}(\beta h)| \cdot
\bigl| \sum_{|k-\nu|\le 2\beta\sigma}
\bigl( \xi(k)\bar\xi(k+h) - \widehat{\chi}(h) \bigr) V_R(k; h)\, \bigr| + o(1) \\
&= \beta\, \bigl(\, \sum_{|h|\le \beta^{-q}} + \sum_{|h|>\beta^{-q}} \,\bigr)\,
|\widehat{g}(\beta h)| \\ &\qquad\qquad\qquad\qquad \times
\bigl|\, \sum_{|k-\nu|\le 2\beta\sigma}
\bigl( \xi(k)\bar\xi(k+h) - \widehat{\chi}(h) \bigr) V_R(k; h) \,\bigr| + o(1)\,.
\end{align*}
For $|h|>\beta^{-q}$, using the crude estimate
\[
\bigl|\, \sum_{|k-\nu|\le 2\beta\sigma}
\bigl( \xi(k)\bar\xi(k+h) - \widehat{\chi}(h) \bigr) V_R(k; h) \,\bigr|
\lesssim \beta\sigma \cdot \sqrt{\sigma} = \beta\sigma^{3/2}\,,
\]
we get
\begin{multline*}
\beta\, \sum_{|h|> \beta^{-q}}
|\widehat{g}(\beta h)| \cdot
\bigl|\, \sum_{|k-\nu|\le 2\beta\sigma}
\bigl( \xi(k)\bar\xi(k+h) - \widehat{\chi}(h) \bigr) V_R(k; h) \,\bigr|
\\
\lesssim \beta^2\sigma^{3/2} \sum_{|h|> \beta^{-q}}
|\widehat{g}(\beta h)|\,,
\end{multline*}
which is $\ll \beta^2\sigma^{3/2}\ep_2(\beta)$ by assumption (d).

Now, we consider the sum over $|h|\le\beta^{-q}$.
Applying summation by parts and using Lemma~\ref{Lemma5D}, we have
\begin{align*}
\Bigl| &\sum_{|k-\nu|\le 2\beta\sigma}
\bigl( \xi(k)\bar\xi(k+h) - \widehat{\chi}(h) \bigr) V_R(k; h) \Bigr|
\\
&\lesssim
\Bigl(\max_{|k-\nu|\le 2\beta\sigma} V_R(k; h)+ \sum_{|k-\nu|\le 2\beta\sigma} |V_R(k; h)-V_R(k-1; h)|\,\Bigr) 
\\&\qquad\qquad\qquad\qquad\qquad\times
\max_{|k-\nu|\le 2\beta}\, \bigl|\,
\sum_{k\le s \le \nu +2\beta\sigma}
\bigl(  \xi(s)\bar\xi(s+h) - \widehat{\chi}(h) \bigr)\,
\bigr|
\\
&\lesssim
\sqrt{\sigma}\, \max_{|k-\nu|\le 2\beta}\, \bigl|\,
\sum_{k\le s \le \nu +2\beta\sigma}
\bigl(  \xi(s)\bar\xi(s+h) - \widehat{\chi}(h) \bigr)
\,\bigr|\,.
\end{align*}
First, we assume that $0\le h \le \beta^{-q}$. Then, by estimate~\eqref{eq:C1},
the maximum on the RHS is $ \lesssim (\nu + 2\beta\sigma)\ep_1(\nu-2\beta\sigma; h)$.
Since $\beta\sigma \ll \sigma/\Delta(\sigma) \lesssim \nu$, the latter expression is
$ \lesssim \, \nu\ep_1(\tfrac12\, \nu; h) $. The application of estimate~\eqref{eq:C1} was legal since,
$  \beta^{-q} \stackrel{(a)}\le H(\tfrac12\, \nu) \,
\stackrel{\beta\sigma\ll\nu}\le\, H(\nu-2\beta\sigma) $.
Thus,
\begin{multline*}
\beta\, \sum_{0\le h \le \beta^{-q}}\,
|\widehat{g}(\beta h)| \cdot
\bigl|\, \sum_{|k-\nu|\le 2\beta\sigma}
\bigl( \xi(k)\bar\xi(k+h) - \widehat{\chi}(h) \bigr) V_R(k; h) \,\bigr|
\\
\lesssim
\beta\sqrt{\sigma}\cdot \nu\, \sum_{0\le h \le \beta^{-q}} \ep_1(\tfrac12\, \nu, h)\,,
\end{multline*}
which is $\ll \beta^2\sigma^{3/2} \ep_2(\beta)$ by assumption (c).

The sum over $-\beta^{-q} \le h \le -1$ needs only a minor modification.
In this case we have
\begin{align*}
\bigl|\, \sum_{k\le s \le \nu +2\beta\sigma}
\bigl(  \xi(s)\bar\xi(s+h) &- \widehat{\chi}(h) \bigr) \,\bigr|\\
&= \bigl|\, \sum_{k-|h|\le s \le \nu +2\beta\sigma-|h|}\,
\bigl(  \xi(s)\bar\xi(s+|h|) - \widehat{\chi}(|h|) \bigr) \,\bigr| \\
&= \bigl|\, \sum_{k\le s \le \nu +2\beta\sigma}\,
\bigl(  \xi(s)\bar\xi(s+|h|) - \widehat{\chi}(|h|) \bigr) \,\bigr| + O(|h|)\,.
\end{align*}
Therefore,
\begin{multline*}
\beta\, \sum_{-\beta^{-q}\le h \le -1}\,
|\widehat{g}(\beta h)| \cdot
\bigl|\, \sum_{|k-\nu|\le 2\beta\sigma}
\bigl( \xi(k)\bar\xi(k+h) - \widehat{\chi}(h) \bigr) V_R(k; h) \,\bigr|
\\
\lesssim
\beta\sqrt{\sigma}\bigl(\, \nu\, \sum_{0\le h' \le \beta^{-q}} \ep_1(\tfrac12\, \nu, h')
+ \beta^{-2q}\,\bigr)\,,
\end{multline*}
and, by assumptions (c) and (b), both terms on the RHS are $\ll \beta^2\sigma^{3/2} \ep_2(\beta)$.

\medskip\noindent\underline{Upper bound of $III$}:
We have
\[
| III | \lesssim \beta\, \bigl(\, \sum_{|h|\le\sqrt\sigma} + \sum_{|h|>\sqrt{\sigma}} \,\bigr)
|\widehat{g}(\beta h)|\cdot |(V_R(h) - V_R(0))|\,.
\]
To estimate the first sum, we apply Lemma~\ref{Lemma10b} and use that the Fourier transform of $g$ decays faster than any negative power. We get
\begin{multline*}
\beta\, \sum_{|h|\le\sqrt\sigma}
|\widehat{g}(\beta h)|\cdot |(V_R(h) - V_R(0))| \lesssim \beta\sqrt{\sigma}\,
\sum_{|h|\le\sqrt\sigma}\,  |\widehat{g}(\beta h)| (1+ h^2\beta) \\
\lesssim \beta\sqrt{\sigma} \cdot \beta^{-2} = \frac{\sqrt{\sigma}}\beta\,.
\end{multline*}
To estimate the second sum, we use the crude bound $V_R(h) \lesssim \beta\sigma^{3/2}$,
which follows from Lemma~\ref{Lemma5C}(i) and Lemma~\ref{Lemma5B}(ii).
Using again that the Fourier transform of $g$ decays faster than any negative power, we get
\begin{multline*}
\beta\, \sum_{|h|>\sqrt\sigma}
|\widehat{g}(\beta h)|\cdot |(V_R(h) - V_R(0))| \lesssim \beta^2\sigma^{3/2}\,
\sum_{|h|>\sqrt\sigma} |\widehat{g}(\beta h)| \\
\lesssim \beta^2 \sigma^{3/2} \cdot \frac1{\beta}\,
\sum_{\ell\ge\beta\sqrt\sigma} \frac1{\ell^3} \lesssim \frac{\sqrt{\sigma}}\beta\,.
\end{multline*}
It remains to recall that
condition (b) guarantees that 
$$
\beta^{-1}\sqrt{\sigma} \ll \beta^2 \sigma^{3/2}\, \ep_2(\beta).
$$
This completes the proof of Proposition~\ref{Lemma5a}.
\end{proof}
\end{document}
